% ============================================================
%  Hands-On 2 — Instalasi Ubuntu pada VirtualBox
% ============================================================

\chapter{Automasi dengan Bash Scripting}
\label{chap:handson-02}

% --- 1. Tugas ---
\section{Tugas}

Jawablah pertanyaan-pertanyaan berikut berdasarkan hasil praktikum Anda.
Sertakan bukti berupa \textit{screenshot} atau penggalan keluaran terminal
yang relevan pada setiap jawaban.

\begin{enumerate}
    \item \textbf{Tujuan dan Output Praktikum}\\[0.3cm]
    \textbf{A. Tujuan Praktikum}
    \begin{enumerate}[nosep]
        \item Menjelaskan peran Sistem Operasi Linux beserta komponen intinya, seperti kernel dan shell.
        \item Menguraikan hierarki sistem berkas Linux dan melakukan navigasi direktori melalui Antarmuka Baris Perintah (CLI).
        \item Menerapkan aturan sintaks perintah Linux, termasuk penggunaan flag dan argumen secara tepat.
        \item Mengoperasikan perintah-perintah dasar untuk manajemen berkas dan direktori (seperti ls, cd, mkdir, rm, cp, mv, touch, dan cat).
        \item Memahami konsep kepemilikan (ownership) dan hak akses (permissions) pada lingkungan berbasis Unix.
        \item Mengonfigurasi hak akses dan kepemilikan menggunakan chmod dan chown, serta memahami penggunaan otorisasi superuser (sudo).
    \end{enumerate}
    \textbf{B. Output Praktikum}
    \begin{enumerate}[nosep]
        \item Struktur direktori dan berkas rekayasa hasil manipulasi (pembuatan, penyalinan, pemindahan) yang dilakukan melalui terminal CLI.
        \item Dokumentasi log modifikasi hak akses berkas menggunakan kombinasi izin read, write, dan execute.
        \item Pemahaman perbedaan eksekusi perintah antara pengguna biasa (\$) dan superuser (\#).
    \end{enumerate}
\end{enumerate}

% --- 2. Kesimpulan ---
\newpage
\section{Kesimpulan}

\noindent\textit{[TODO: Tuliskan kesimpulan dari praktikum yang telah dilaksanakan.
Kesimpulan hendaknya menjawab tujuan-tujuan pada modul praktikum, minimal 2--3 paragraf.]}

% --- 3. Lampiran: Bukti Penggunaan AI ---
\newpage
\section{Lampiran: Bukti Penggunaan AI}
\label{sec:h02-lampiran-ai}

Sesuai dengan kebijakan penggunaan kecerdasan buatan (AI) dalam kegiatan akademik,
praktikan \textbf{wajib} mencantumkan setiap penggunaan alat bantu AI yang dilakukan
selama pengerjaan hands-on ini. Ketidakjujuran dalam pelaporan penggunaan AI dianggap
sebagai pelanggaran integritas akademik.

\subsection{Identitas Penggunaan AI}

\begin{table}[h]
\caption{Informasi Alat AI yang Digunakan}
\label{tab:h02-ai-info}
\centering
\begin{tabulary}{\textwidth}{|L|L|}
\hline
\textbf{Keterangan}   & \textbf{Isian} \\ \hline
Nama alat AI          & \textit{(contoh: ChatGPT, GitHub Copilot, Gemini, dst.)} \\ \hline
Versi / Model         & \textit{(contoh: GPT-4o, Gemini 1.5 Pro, dst.)} \\ \hline
Tanggal penggunaan    & \textit{TODO: isi tanggal} \\ \hline
Tujuan penggunaan     & \textit{TODO: jelaskan secara singkat untuk apa AI digunakan} \\ \hline
\end{tabulary}
\end{table}

\subsection{Tangkapan Layar Percakapan AI}

Sertakan tangkapan layar (\textit{screenshot}) percakapan dengan AI sebagai
bukti otentik. Setiap gambar harus memperlihatkan \textbf{prompt} yang dikirimkan
dan \textbf{respons} yang diterima secara jelas.

\begin{figure}[h]
    \centering
    % TODO: ganti dengan screenshot percakapan AI Anda
    \includegraphics[width=0.85\textwidth]{Figure/ifitera-header.png}
    \caption{Tangkapan Layar Percakapan AI --- Sesi 1 (TODO: ganti caption)}
    \label{fig:h02-ai-screenshot-1}
\end{figure}
